\documentclass[journal]{IEEEtran}

\usepackage{amsmath,amssymb}
\usepackage{graphicx}
\usepackage{booktabs}
\usepackage{multirow}
\usepackage{array}
\usepackage{cite}
\usepackage{url}
\usepackage{dblfloatfix}
\usepackage{float}

\title{An Integrated IoT Framework for Real-Time Environmental Monitoring, Forecasting, and Anomaly Detection}

\author{
Pratik~Kumar,
Nikunj~Goydani,
and~Vetriveeran~Rajamanai
\thanks{
All authors are with Vellore Institute of Technology, Vellore, India.
Pratik Kumar (e-mail: pratik.kumar2024a@vitstudent.ac.in);
Nikunj Goydani (e-mail: nikunj.goydani2024@vitstudent.ac.in);
Vetriveeran Rajamanai (e-mail: vetriveeran.r@vit.ac.in).
Vetriveeran Rajamanai is the corresponding author.
}
}

\begin{document}

\maketitle

% -------------------------------------------------
% Abstract
% -------------------------------------------------
\begin{abstract}
Environmental monitoring using resource-constrained Internet of Things (IoT)
nodes is limited by the availability of analog-to-digital converter (ADC)
channels, while existing studies often address multi-parameter sensing, cloud
monitoring, forecasting, and anomaly detection as separate components. This
study presents an integrated and cost-efficient environmental monitoring
framework that uses a CD4051 analog multiplexer with an ESP8266 to acquire
temperature, humidity, gas concentration, light intensity, and soil-moisture
measurements through a single ADC channel, together with wireless data
transmission, backend processing, and machine-learning analytics. The system
was evaluated using five days of laboratory measurements collected at 10-s
intervals. Independent Long Short-Term Memory (LSTM) models were developed
for parameter-wise short-term forecasting, and Isolation Forest was employed
for unsupervised anomaly detection using only recorded sensor data. The system
provided reliable continuous data transmission with no observable
channel-to-channel crosstalk during multiplexed acquisition. The analytical
evaluation achieved 94\% agreement across 9,046 recorded observations, based
on comparison with reference measurements obtained from existing laboratory
equipment. Temperature recorded a mean of 22.89~$^\circ$C
(19.70--26.30~$^\circ$C), humidity 49.82\% (38.50--60.50\%), gas level
77.68 (61--94), and light intensity 533.73 (373--900). Strong correlations
were observed between temperature and humidity ($-0.89$), temperature and
light intensity ($+0.85$), and humidity and light intensity ($-0.87$).
Overall, the proposed framework demonstrates the feasibility of integrating
ADC-efficient multi-parameter sensing with forecasting and anomaly detection
in a single low-resource IoT node, providing a cost-efficient alternative to
architectures requiring additional ADC resources or separately deployed
analytical components.
\end{abstract}

% -------------------------------------------------
% Keywords
% -------------------------------------------------
\begin{IEEEkeywords}
Environmental Monitoring, Internet of Things (IoT), Analog Multiplexing,
ESP8266, CD4051, Multi-Parameter Sensing, LSTM Forecasting,
Anomaly Detection, Isolation Forest, Predictive Analytics
\end{IEEEkeywords}

% =================================================
% I. INTRODUCTION
% =================================================
\section{Introduction}

In the last few years, the importance of environmental monitoring has
increased significantly as regards the smart building sector, the precision
agriculture industry, industrial safety and laboratory quality assurance
\cite{ref1}. The increased use of low-cost microcontrollers and their ability
to connect to networks has made it possible to create devices that can use
these microcontrollers to form distributed sensor networks capable of
measuring temperature, humidity, air quality, light, and soil conditions in
real-time, as well as transmitting that data to cloud or edge analytics
platforms \cite{ref2}. The research in this area covers the integration of
hardware, the development of communication protocols, the storage of data,
the creation of data visualizations, and how machine learning (ML) can be
applied for detection of anomalies and predictive analytics \cite{ref3}.

One of the most significant challenges with low-cost Internet of Things
(IoT) nodes \cite{ref4} is the lack of analog-to-digital converter (ADC)
resources. Many of the commonly used microcontrollers feature only one or a
few ADC channels, and as a result, they have a limited number of analog
sensors that can be connected without using either an external expansion
method or multiple devices. Historically, analog multiplexers (MX) have been
used in embedded systems to allow multiple analog sources to share the same
ADC \cite{ref5}.

Current research is focused on the incorporation of a variety of
environmental factors into one single monitoring system. This process
includes data collection, data storage, data visualization, and using
machine learning processes for prediction of data and detection of
anomalies \cite{ref6}. The significance of the research is emphasized by
the observation that models have been assessed using real laboratory and
actual data. The paper describes the development and assessment of the
environment monitoring device which consists of three components: a
hardware acquisition component which makes use of a CD4051 analog
multiplexer along with an ESP8266 microcontroller \cite{ref7}. The device
can measure different parameters such as gas, light and moisture, digital
temperature and humidity by means of one ADC. The second component includes
a backend and a web-based dashboard which are used for storing and
visualizing the data \cite{ref8}. The last component is the machine
learning pipeline which performs time-series forecasting with the use of
LSTM algorithms for data collected in laboratory environments \cite{ref9}.

This study is positioned as a research-oriented investigation into the
co-design of hardware, software, and machine-learning components,
emphasizing reproducible implementation strategies, structured data
handling practices, and systematic evaluation methodologies. The principal
contributions include the development of a documented interfacing
methodology for integrating a CD4051 MX with an ESP8266 microcontroller
\cite{ref10}, design approaches for constructing an end-to-end
Environmental Monitoring Device (EMD) \cite{ref11}, the generation of a
multi-parameter temporal environmental dataset collected under controlled
laboratory conditions, and the realization of a structured machine-learning
framework enabling prediction and anomaly detection using exclusively real
laboratory measurements \cite{ref12}.
%\section*{Abbreviations}

The abbreviations used throughout this manuscript are summarized in
Table~\ref{tab:abbreviations}.

\begin{table}[!t]
    \centering
    \caption{Abbreviations used in this manuscript.}
    \label{tab:abbreviations}
    \renewcommand{\arraystretch}{1.08}
    \setlength{\tabcolsep}{5pt}

    \begin{tabular}{ll}
        \toprule
        \textbf{Abbreviation} & \textbf{Definition} \\
        \midrule

        ADC  & Analog-to-Digital Converter \\
        API  & Application Programming Interface \\
        DAC  & Digital-to-Analog Converter \\
        CoAP & Constrained Application Protocol \\
        DHT22 & Digital temperature/humidity sensor (model) \\
        EMD  & Environmental Monitoring Device \\
        GPIO & General Purpose Input/Output \\
        I2C  & Inter-Integrated Circuit (serial bus) \\
        IoT  & Internet of Things \\
        LDR  & Light-Dependent Resistor \\
        MAE  & Mean Absolute Error \\
        MCU  & Microcontroller Unit \\
        ML   & Machine Learning \\
        MQTT & Message Queuing Telemetry Transport \\
        MX   & Multiplexer \\
        POST & HTTP request method \\
        REST & Representational State Transfer \\
        RH   & Relative Humidity \\
        RMSE & Root Mean Square Error \\
        SCL  & Serial Clock (I2C line) \\
        LPG  & Liquefied Petroleum Gas \\
        LSTM & Long Short-Term Memory \\

        \bottomrule
    \end{tabular}
\end{table}


% -------------------------------------------------
% Research Gap
% -------------------------------------------------
\subsection{Research Gap and Problem Statement}

The widespread use of environmentally monitored Internet of Things systems
is being hampered by the availability of low-cost microcontroller hardware
for practical use in an Multi-Parameter Sensing Architecture (MPSA). The
limitations in these hardware products, due to a lack of ADC capabilities,
hinder the ability to integrate multiple analogue sensors. Most current
studies on sensing, backend data management, or ML analysis examine each
type of data as independent from one another rather than as a unified pair,
concluding with multiple extreme solutions that have limited scalability or
reproducibility. Throughout the majority of ML-based monitoring studies,
there are instances where either synthetic data or augmented data is used
for producing results, creating difficulty in creating reliable testing
results in actual laboratory and industrial environments.

% -------------------------------------------------
% Proposed Framework
% -------------------------------------------------
\subsection{Proposed Framework Overview}

In an effort to improve upon identified shortcomings, the proposed work
introduces a comprehensive environmental monitoring platform which
incorporates the use of multiplexed analog sensors, scalable cloud data
storage and visualization capabilities as well as machine learning
analytics all part of an overarching architectural framework. In order to
achieve improved analog sensing capabilities for use with the ESP8266
microcontroller, the utilization of a CD4051 analog multiplexer will be
essential to allow for efficient multi-parameter data collection using
multiplexed analog sensors. The proposed system includes a Flask-based
backend for data management, storage and visualization, as well as a
React-based dashboard for live monitoring of environmental conditions.
Additionally, the proposed framework also integrates predictive and
anomaly detection models, specifically long short-term memory (LSTM) and
isolation forest algorithms, that have only been trained on representative
laboratory-assessed observations in order to ensure the successful
implementation of the proposed monitoring solution in a practical
application and operationally reliable manner.

% -------------------------------------------------
% Contributions
% -------------------------------------------------
\subsection{Contributions of the Study}

The principal contributions of this study are as follows:

\begin{itemize}
    \item A flexible, multi-sensor interface that allows for multi-parameter
    environmental monitoring utilizing a single ADC microcontroller.

    \item Complete integration of a total environmental monitoring system,
    from data collection to backend storage, to real-time display.

    \item Creation of a natural laboratory environmental dataset that can
    be utilized in the evaluation of the system, without reliance on
    synthetic data.

    \item Development of a fully functional ML analytics framework consisting
    of LSTM time-series forecast and isolation forest anomaly detection.

    \item Documentation of the hardware-software co-design processes in
    order to provide for reproducibility and to allow for additional
    research to build on existing work.
\end{itemize}

% =================================================
% II. MATERIALS AND METHODS
% =================================================
\section{Materials and Methods}

This section describes the system architecture of the Environmental
Monitoring Device (EMD). The hardware components involved in the
multi-parameter sensing system, including ADC multiplexing, backend
processing and data management, and real-time visualization, are discussed.
The workflow of the machine-learning algorithms used for prediction and
anomaly detection is also presented.

In general, the proposed architecture consists of six integrated layers.
First, the edge sensing and data acquisition layer employs ESP8266-based
sensor nodes for acquiring multiple environmental parameters. Second, the
hardware design and integration layer uses a CD4051 analog multiplexer to
enable multiple analog sensors to share a single ADC channel of the
ESP8266. Third, the communication and backend data management layer is
implemented using Flask for data ingestion and access. Fourth, the database
and service layer provides persistent storage for the collected time-series data. Fifth,
the visualization and user interaction layer provides authentication,
configuration management, and a React-based frontend dashboard for
real-time monitoring and analytics. Finally, the machine-learning and
analytics layer processes the collected sensor data for forecasting and
anomaly detection.

The architecture enables the ESP8266 to acquire multiple analog
environmental parameters through a single ADC channel using CD4051
multiplexing, thereby addressing the hardware constraints of low-cost IoT
nodes. The acquired measurements are transmitted to the backend for storage
and subsequent visualization and machine-learning analysis. This integrated
data pathway supports continuous monitoring while providing the foundation
for forecasting and anomaly detection. The overall architecture integrates
sensor data acquisition, multiplexed ADC conversion, backend processing,
time-series data storage, real-time visualization, and machine-learning-based
analysis into a unified environmental monitoring framework. The overall
system architecture and data flow of the proposed EMD are presented in
Fig.~\ref{fig:architecture}.

% -------------------------------------------------
% Figure 1
% -------------------------------------------------
\begin{figure*}[!t]
    \centering
    \includegraphics[width=\textwidth]{fig1.png}
    \caption{Overall architecture of the proposed Environmental Monitoring
    Device (EMD), integrating edge sensing, ADC multiplexing, backend data
    management, time-series storage, visualization, and machine-learning-based
    forecasting and anomaly detection.}
    \label{fig:architecture}
\end{figure*}
% =================================================
% A. EDGE SENSING AND DATA ACQUISITION
% =================================================
\subsection{Edge Sensing and Data Acquisition}

The proposed node utilizes an ESP8266MOD microcontroller, which provides a
single analog-to-digital converter (ADC) input through the A0 pin. The analog
signals from the MQ-6 gas sensor, light-dependent resistor (LDR), and
soil-moisture sensor are sequentially connected to the ADC through the CD4051
analog multiplexer. This arrangement enables multiple analog sensors to share
the single ADC input of the ESP8266.

The following sensors and interface components are integrated with the
proposed node:

\begin{itemize}
    \item \textbf{DHT22:} The DHT22 is a digital temperature and humidity
    sensor connected to GPIO2 of the ESP8266 through a 10~k$\Omega$ resistor.
    Temperature and humidity measurements are sampled at 10-s intervals.

    \item \textbf{MQ-6:} The MQ-6 gas sensor is powered by 5~V and is used
    for gas concentration measurement, including LPG-related detection.
    Its analog output is connected to channel 0 of the CD4051. After
    selecting the corresponding channel, the analog signal is converted
    into a digital value through the ESP8266 A0 input.

    \item \textbf{LDR:} The light-dependent resistor is connected to
    channel 1 of the CD4051. Similar to the MQ-6, its analog signal is
    routed to the ESP8266 ADC through the multiplexer. A higher ADC value
    indicates greater light intensity.

    \item \textbf{Soil-Moisture Sensor:} The soil-moisture sensor is
    connected to channel 2 of the CD4051. Its analog reading is acquired
    through the ESP8266 ADC, where a higher value indicates wetter soil.

    \item \textbf{CD4051:} The CD4051 is an 8-channel analog multiplexer
    controlled through the select pins S0, S1, and S2. These control signals
    are connected to GPIO12 (D6), GPIO13 (D7), and GPIO14 (D5),
    respectively. The common output Z is connected to the ESP8266 A0 input.
    Channels 0, 1, and 2 are used for the MQ-6, LDR, and soil-moisture
    sensors, respectively, while channels 3--7 are reserved for future
    expansion.
\end{itemize}

During operation, the firmware selects the CD4051 channels sequentially.
After each channel selection, a short stabilization delay of approximately
1--2~ms is introduced before the analog value is acquired through the A0
input. Every 10~s, the measurements from all five environmental parameters
are combined into a single JSON message and transmitted to the backend using
an HTTP POST request. In the event of network connectivity problems or
sensor-related failures, an error message is displayed on the LCD and the
transmission is retried.


% =================================================
% B. HARDWARE DESIGN AND INTEGRATION
% =================================================
\subsection{Hardware Design and Integration}

The prototype Environmental Monitoring Device (EMD) integrates sensing,
signal multiplexing, processing, local display, power management, and wireless
communication within a single IoT node. The physical implementation of the
prototype is shown in Fig.~\ref{fig:hardware}, while the corresponding
electrical and communication architecture is presented in
Fig.~\ref{fig:communication}.

% -------------------------------------------------
% Figure 2
% -------------------------------------------------
\begin{figure}[!t]
    \centering
    \includegraphics[width=\columnwidth]{fig2.png}
    \caption{Hardware prototype of the proposed Environmental Monitoring
    Device (EMD), showing the ESP8266 microcontroller, CD4051 analog
    multiplexer, DHT22 temperature and humidity sensor, MQ-6 gas sensor,
    LDR, soil-moisture sensor, JHD162A LCD module, and power module.}
    \label{fig:hardware}
\end{figure}

The ESP8266MOD acts as the main control unit, performing sensor acquisition,
data processing, peripheral control, and Wi-Fi communication. Since the
ESP8266 provides only one analog input, the CD4051 analog multiplexer is used
to connect the three analog sensors, namely the MQ-6 gas sensor, LDR, and
soil-moisture sensor, sequentially to the ADC. The ESP8266 selects the
required CD4051 channel by applying the appropriate digital control signals
to the multiplexer and subsequently acquires the corresponding sensor signal
through its ADC. In contrast, the DHT22 temperature and humidity sensor is
connected directly to the ESP8266 through its digital interface and therefore
does not require the multiplexer.

The ESP8266 operates in a continuous acquisition cycle. It first obtains
temperature and humidity measurements from the DHT22 and subsequently
acquires measurements from the analog sensors through the CD4051. The
multiplexer enables rapid switching between the analog inputs; however, a
short stabilization period is provided after each channel transition before
the ADC conversion is performed. The processed measurements are then
transmitted from the microcontroller to the network router as JSON messages,
providing a simple data representation for backend processing.

Figure~\ref{fig:hardware} illustrates the physical arrangement and wiring of
the prototype, including the placement and interconnection of the sensing,
processing, display, and power components. The JHD162A 16$\times$2 LCD is
connected to the ESP8266 through the I$^2$C interface for local display of
sensor readings and system status. The processed measurements are transmitted
through the ESP8266 Wi-Fi interface to the backend system, thereby providing
wireless communication between the sensing node and the data-management
layer.

During operation, the ESP8266 sequentially controls the MQ-6, LDR, and
soil-moisture channels through the CD4051, allowing sufficient stabilization
time before each ADC conversion. Temperature and humidity measurements are
obtained directly from the DHT22 through its digital interface. Thus,
Fig.~\ref{fig:hardware} presents the physical implementation of the proposed
EMD, whereas Fig.~\ref{fig:communication} illustrates the relationships
among the sensor signals, multiplexer control, power supply, local display,
and wireless communication.


% -------------------------------------------------
% Figure 3
% -------------------------------------------------
\begin{figure*}[!t]
    \centering
    \includegraphics[width=\textwidth]{fig3.png}
    \caption{Electrical and communication architecture of the proposed
    Environmental Monitoring Device (EMD), illustrating sensor interfaces,
    CD4051 channel selection, ESP8266 processing, local display, power
    connections, and wireless communication with the backend system.}
    \label{fig:communication}
\end{figure*}
% =================================================
% C. COMMUNICATION AND BACKEND DATA MANAGEMENT
% =================================================
\subsection{Communication and Backend Data Management}

The backend system was developed using the Flask framework and provides
REST-based services for data ingestion, retrieval, statistical analysis,
system monitoring, and machine-learning inference. The principal API
endpoints are described as follows:

\begin{itemize}

    \item \textbf{Sensor Data Ingestion:} A POST request from the ESP8266 is
    sent to the \texttt{/api/sensor-data} endpoint. The endpoint stores
    timestamped temperature and humidity measurements, \texttt{gas\_analog},
    \texttt{light\_level}, and \texttt{soil\_moisture} values in the
    database.

    \item \textbf{Latest Measurements:} A GET request to
    \texttt{/api/latest} retrieves the most recent measurements of all five
    environmental parameters for display on the dashboard.

    \item \textbf{Statistical Data:} A GET request to
    \texttt{/api/stats} retrieves statistical information, including the
    minimum, maximum, and average values, for a specified time period.

    \item \textbf{System Health:} A GET request to
    \texttt{/api/health} provides the operational status of the backend and
    database services.

    \item \textbf{Prediction Services:} The endpoints
    \texttt{/predict/temperature}, \texttt{/predict/humidity},
    \texttt{/predict/gas}, \texttt{/predict/light}, and
    \texttt{/predict/soil} provide LSTM-based predictions for the selected
    environmental parameter over the requested prediction horizon, such as
    the next 1--168~h.

\end{itemize}


% =================================================
% D. DATABASE AND SERVICE LAYER
% =================================================
\subsection{Database and Service Layer}

The database and service layer provides persistent storage and supporting
services for the environmental monitoring framework. The collected sensor
measurements are stored as timestamped records, enabling efficient retrieval
of historical observations for visualization, statistical analysis, and
machine-learning applications.

For cloud deployment, the database layer supports scalable cloud-based
storage, allowing the collected environmental data to remain accessible to
the backend services and analytical modules. The service layer manages the
interaction between the stored data and the application components,
providing structured access to sensor measurements, statistical information,
system status, and machine-learning outputs. This separation of data
persistence and service functionality enables the monitoring framework to
maintain continuous data availability while supporting subsequent analysis
and visualization.


% =================================================
% E. VISUALIZATION AND USER INTERACTION
% =================================================
\subsection{Visualization and User Interaction}

Input validation and error handling are implemented to ensure reliable
interaction between the application and the underlying data services. The
visualization layer maintains the same data structure for both local
SQLite-based deployment and cloud-based PostgreSQL deployment, enabling
consistent data handling across deployment environments.

\begin{itemize}

    \item \textbf{Dashboard:} The dashboard is implemented as a single-page
    web application using React with Vite, together with Material UI and
    Plotly libraries. It provides real-time monitoring, historical data
    visualization, prediction display, reporting, and system-status
    information.

    \item \textbf{Real-Time and Archived Data Visualization:} Users can
    visualize the five environmental parameters with an automatic-update
    option. A selectable time range is provided, together with zooming and
    panning functions for detailed examination of the recorded data.

    \item \textbf{LSTM Predictions:} The dashboard provides graphical
    visualization of LSTM predictions for the selected environmental
    parameters over a specified prediction horizon, such as the next 24~h.
    The prediction values are internally normalized, while the corresponding
    original-scale values are displayed when the user hovers over the
    plotted data.

    \item \textbf{PDF Export:} A monitoring report can be generated and
    downloaded in PDF format through the Reports page or by invoking the
    \texttt{POST /api/reports/generate} endpoint with the PDF format.

    \item \textbf{System Status:} The dashboard displays the date and time
    of the latest update together with the operational state of the backend
    and network connection.

    \item \textbf{Alert Settings:} An Alerts page allows users to configure
    threshold values for temperature (minimum and maximum), humidity
    (minimum and maximum), light level/LDR (minimum and maximum), soil
    moisture (minimum and maximum), and gas concentration (maximum). The
    configured thresholds are stored in the browser using
    \texttt{localStorage}. The page also displays anomalies detected by the
    backend through the \texttt{GET /api/anomalies} endpoint for the previous
    24~h, including timestamps and corresponding sensor values.

\end{itemize}


% =================================================
% F. MACHINE LEARNING AND ANALYTICS
% =================================================
\subsection{Machine Learning and Analytics}

The machine-learning pipeline processes only recorded sensor data, such as
the measurements obtained during the five-day laboratory deployment.
Synthetic or artificially generated time-series data are not included in
the analysis. The major stages of the machine-learning pipeline are
described below.

\begin{itemize}

    \item \textbf{Data:} The time-series measurements of temperature,
    humidity, \texttt{gas\_analog}, \texttt{light\_level}, and
    \texttt{soil\_moisture}, together with their timestamps, are retrieved
    from the database. Synthetic or generated time-series data are excluded
    from the training data.

    \item \textbf{Preprocessing:} Rolling means and change-based features
    are generated for the environmental parameters. Additional temporal
    features, such as time of day and day of week, can optionally be
    incorporated. The data are normalized using methods such as
    \texttt{MinMaxScaler}. Missing observations are handled before model
    training, for example, by replacing missing values with the corresponding
    column mean.

    \item \textbf{Training:} Independent LSTM models are developed for each
    of the five environmental parameters. The models are trained using
    sequences of previous observations to predict future values.
    Isolation Forest is employed to identify anomalous observations from
    the same recorded dataset.

    \item \textbf{Evaluation:} The models are evaluated using a hold-out or
    time-based data split. Regression performance is assessed using metrics
    such as mean absolute error (MAE) and root mean square error (RMSE),
    whereas anomaly detection performance is evaluated using measures such
    as precision and recall.

    \item \textbf{Deployment:} The trained models and corresponding
    normalization scalers are loaded by the Flask backend. Cloud-based
    inference is performed through the prediction endpoints, enabling the
    generated forecasts to be accessed by the frontend monitoring system.

\end{itemize}

The collected time-series data are therefore processed using rolling and
change-based features within the machine-learning pipeline. Independent
LSTM models are developed for each environmental parameter, while anomalies
are detected using Isolation Forest. The prediction inference process begins
with the backend prediction services and converts the recorded sensor data
into forecasting and anomaly-related information for integration with the
monitoring system. The complete machine-learning workflow is illustrated in
Fig.~\ref{fig:ml_pipeline}.


% -------------------------------------------------
% Figure 4
% -------------------------------------------------
\begin{figure*}[!t]
    \centering
    \includegraphics[width=\textwidth]{fig4.png}
    \caption{Machine-learning pipeline and analytics layer of the proposed
    EMD, showing the sequential stages of recorded sensor-data retrieval,
    preprocessing, model training, evaluation, and deployment. The pipeline
    uses LSTM models for parameter-wise forecasting and Isolation Forest for
    anomaly detection, producing forecasts and anomaly-related outputs for
    integration with the monitoring system.}
    \label{fig:ml_pipeline}
\end{figure*}
% =================================================
% III. RESULTS
% =================================================
\section{Results}

The experimental evaluation was conducted over five days of continuous
operation in the laboratory. Sensor data were collected at 10-s intervals
and subsequently processed by the backend system. To avoid data leakage, a
strict temporal split was applied to preserve the chronological order of
the recorded data between training and testing datasets. Quality-of-service
checks were performed to verify that the time-series data maintained
consistent duration and sampling intervals before machine-learning-based
prediction. This approach provides consistent conditions for evaluating
LSTM forecasting and Isolation Forest anomaly detection across multiple
runs.

\subsection{Sensor Data Acquisition}

Sensor data were continuously collected for five environmental parameters:
temperature, humidity, \texttt{gas\_analog}, \texttt{light\_level}, and
\texttt{soil\_moisture}. The DHT22 provided temperature and humidity
measurements, while the analog parameters, namely gas, light, and soil
moisture, were acquired through CD4051 multiplexed switching with a
settling time of approximately 1--2~ms after each channel selection.
Initial signal-stability tests indicated only a small amount of noise in
the recorded data. Therefore, the collected measurements were directly
used for subsequent training and evaluation.

\subsection{Backend and User Interface}

The Flask-based backend APIs received and processed the sensor measurements,
while the React-based interface provided visualization of both real-time
and historical data for all five parameters. The interface also provided
prediction and report-export functionalities. The responses from the
aggregation and prediction services exhibited acceptable latency during
laboratory testing, enabling interactive monitoring and smooth updating of
the displayed visualizations.

\subsection{Machine Learning Results}

LSTM models were trained independently for the five target parameters using
the data collected during the five-day laboratory deployment. The models
learned temporal dependencies from previous observations to generate
short-term predictions. An Isolation Forest model was trained using the
same recorded sensor data for anomaly detection. Model evaluation was
performed using previously unseen temporal segments of the dataset, thereby
maintaining chronological separation between training and evaluation data.

The anomaly detection results identified significant deviations in
parameters such as gas level and temperature, while the LSTM predictions
followed the short-term trends observed in the recorded measurements.

% -------------------------------------------------
% Figure 5
% -------------------------------------------------
\begin{figure}[!t]
    \centering
    \includegraphics[width=\columnwidth]{fig5.png}
    \caption{Real-time visualization interface of the proposed EMD,
    displaying temperature, relative humidity, gas level, light intensity,
    and soil-moisture measurements together with the latest data-update
    time and 24-h statistical summaries.}
    \label{fig:dashboard}
\end{figure}

The React-based Environmental Monitoring Dashboard is shown in
Fig.~\ref{fig:dashboard}. The dashboard was developed and tested in the
laboratory for real-time monitoring of environmental conditions. During the
laboratory test, the device measured a temperature of 22.2~$^\circ$C and a
relative humidity of 55.1\%. The gas-level reading was 70/1024, indicating
a low sensor response, while the light-level reading was 374/1024,
representing a moderate level of surrounding light. The soil-moisture
measurement was 500/1024, representing an intermediate moisture level.
These measurements demonstrate the ability of the proposed device to
capture, process, and display multiple environmental parameters in
real time.

\subsection{Report Generation and Statistical Analysis}

The report-generation feature was successfully implemented and tested under
laboratory conditions. Users can specify a date range to generate a report
containing statistical summaries of the monitored environmental parameters.
As shown in Fig.~\ref{fig:report}, the report interface provides the
observation count, mean, median, standard deviation, maximum, and minimum
values for each parameter.

The recorded temperature had a mean of 22.89~$^\circ$C, with values ranging
from 19.70~$^\circ$C to 26.30~$^\circ$C. Humidity had a mean of 49.82\%,
with values ranging from 38.50\% to 60.50\%. The gas-level readings had a
mean of 77.68, ranging from 61 to 94, whereas the light-intensity readings
had a mean of 533.73, ranging from 373 to 900. The soil-moisture value
remained at 500 throughout the recorded period, resulting in a standard
deviation of 0.00 and indicating no variation in the recorded soil-moisture
measurements.

The generated report also incorporates analytical visualizations to provide
a comprehensive representation of the collected sensor data. The
correlation matrix is presented in Fig.~\ref{fig:correlation}. The complete
generated report, including the statistical summary and graphical
representations, is presented in Fig.~\ref{fig:appendix_report}.
These results demonstrate the capability of the reporting module to convert
the collected laboratory sensor data into a structured and interpretable
report for further analysis and evaluation.

% -------------------------------------------------
% Figure 6
% -------------------------------------------------
\begin{figure}[!t]
    \centering
    \includegraphics[width=\columnwidth]{fig6.png}
    \caption{Report-generation interface of the proposed EMD, enabling
    selection of a specific date range and generation of a statistical
    report containing observation count, mean, median, standard deviation,
    minimum, and maximum values for the monitored environmental parameters,
    with an option to export the results as a PDF.}
    \label{fig:report}
\end{figure}

\subsection{Correlation Analysis}

The correlation matrix was used to identify the linear relationships among
the monitored environmental parameters. Pearson correlation coefficients
ranging from $-1$ to $+1$ were calculated, where values closer to
$\pm1$ indicate stronger linear relationships and values close to zero
indicate weaker relationships.

As shown in Fig.~\ref{fig:correlation}, temperature and humidity exhibit a
strong negative correlation of $-0.89$, indicating that humidity generally
decreased as temperature increased within the recorded laboratory data.
Similarly, temperature and light intensity show a strong positive
correlation of $+0.85$, indicating that higher light levels were generally
associated with higher temperatures. A strong negative correlation of
$-0.87$ was also observed between humidity and light intensity, indicating
an inverse relationship between these parameters.

In comparison, gas level and soil moisture show relatively weak
correlations with the other monitored parameters, indicating limited linear
association within the recorded laboratory dataset. Overall, the strongest
relationships were observed among temperature, humidity, and light
intensity, whereas gas level and soil moisture exhibited comparatively
weaker linear associations.

% -------------------------------------------------
% Figure 7
% -------------------------------------------------
\begin{figure}[!t]
    \centering
    \includegraphics[width=\columnwidth]{fig7.png}
    \caption{Correlation matrix of the monitored environmental parameters,
    showing strong negative correlation between temperature and humidity
    ($-0.89$), strong positive correlation between temperature and light
    intensity ($+0.85$), and strong negative correlation between humidity
    and light intensity ($-0.87$).}
    \label{fig:correlation}
\end{figure}

\subsection{Exported Monitoring Report}

The exported environmental monitoring report, shown in
Fig.~\ref{fig:appendix_report}, contains statistical summaries, parameter
ranges, and correlation analysis for the selected monitoring period. The
combination of these analytical features within a single report provides a
structured representation of the recorded environmental measurements and
supports subsequent analysis and evaluation.

As presented in Fig.~\ref{fig:appendix_report}, the exported report covers
the monitoring period of 9--12 January 2026 and summarizes the recorded
sensor statistics, parameter ranges, and correlation relationships, providing
an overall view of the monitored environmental conditions and their
interactions during the selected period.

% -------------------------------------------------
% Figure 8: Exported monitoring report
% -------------------------------------------------
\begin{figure*}[!t]
    \centering
    \includegraphics[width=0.60\textwidth]{figA1.1.png}
        \caption{Exported environmental monitoring report for 9--12 January 2026,
    summarizing the recorded sensor statistics, parameter ranges, and
    correlation relationships during the selected monitoring period.}
    \label{fig:appendix_report}
\end{figure*}

% -------------------------------------------------
% Additional Results: CD4051/ADC Testing and Alerts
% -------------------------------------------------

\subsection{CD4051 and ADC Testing}

No channel-to-channel bleeding or crosstalk was observed while recording the
collected analog channels. Sequential S0/S1/S2 control signals and the applied
settling delays were sufficient to provide reliable multi-channel recording.

\subsection{Alert Configuration and Anomaly Monitoring}

The Alert section was developed with the capability to configure thresholds
for all five measured parameters: temperature, humidity, light intensity,
soil moisture, and gas level. The configured threshold values are stored in
the browser to ensure persistence across working sessions.

In addition to threshold monitoring, the Alert page provides users with a list
of anomalies detected during the previous 24~h using the Isolation Forest
anomaly-detection model. Each anomaly entry contains its timestamp and the
corresponding sensor-reading value, enabling users to identify the environmental
conditions associated with abnormal observations and to remain informed about
the outputs of the machine-learning-based detection model. The alert
configuration interface, threshold configuration, and detected-anomaly
display are presented in Fig.~\ref{fig:appendix_alerts}, where threshold
limits can be defined for temperature, humidity, light intensity, soil
moisture, and gas level. These user-defined limits enable the system to
identify parameter values outside the acceptable range and support timely
detection of abnormal environmental conditions.

% -------------------------------------------------
% Figure 9: Alert configuration interface
% -------------------------------------------------
\begin{figure*}[!t]
    \centering
    \includegraphics[width=\textwidth]{figA2.png}
    \caption{Alert configuration interface of the proposed EMD, providing
    user-defined threshold settings for temperature, humidity, light
    intensity, soil moisture, and gas level to support detection of abnormal
    environmental conditions.}
    \label{fig:appendix_alerts}
\end{figure*}

\subsection{Overall System Performance}

The Environmental Monitoring Device (EMD) successfully demonstrated reliable
data acquisition through the combined CD4051--ESP8266 interface, enabling the
collection of measurements from multiple sensor channels. In addition to the
data-acquisition capability, the EMD integrates backend data processing and
API-based communication, allowing the collected measurements to be transferred
and stored for subsequent use.

The integration of machine-learning methods further extends the analytical
capabilities of the EMD by enabling the development and evaluation of models
using the data collected during system operation. The corresponding
analytical outputs obtained from the recorded measurements are presented in
Figs.~\ref{fig:appendix_ml1}--\ref{fig:appendix_ml3}.

\subsection{Time-Series Observations of the Monitored Parameters}

Figure~\ref{fig:appendix_ml1} presents the time-series variation of gas level
during the monitoring period. The recorded gas-level measurements exhibit
short-term fluctuations around the observed operating range, providing a
temporal representation of gas-level variations in the monitored environment.

% -------------------------------------------------
% Figure 10: Gas-level time series
% -------------------------------------------------
\begin{figure*}[!t]
    \centering
    \includegraphics[width=\textwidth]{figA3.jpg}
    \caption{Time-series variation of gas level during the monitoring period,
    showing the recorded gas-level measurements over time.}
    \label{fig:appendix_ml1}
\end{figure*}

The corresponding temperature and humidity observations are shown in
Fig.~\ref{fig:appendix_ml2}. The two parameters follow an inverse trend over
the recorded observations, which is consistent with the strong negative
correlation of $-0.89$ reported in Fig.~\ref{fig:correlation}.

% -------------------------------------------------
% Figure 11: Temperature and humidity time series
% -------------------------------------------------
\begin{figure*}[!t]
    \centering
    \includegraphics[width=\textwidth]{figA4.jpg}
    \caption{Time-series variation of temperature and humidity during the
    monitoring period, showing their inverse trend over the recorded
    environmental observations.}
    \label{fig:appendix_ml2}
\end{figure*}

Figure~\ref{fig:appendix_ml3} presents the light-level (LDR) and
soil-moisture analog readings recorded during the five-day laboratory
deployment. Both analog channels were multiplexed through the CD4051 and
acquired using the single ADC channel of the ESP8266, confirming that the
multiplexed acquisition scheme described in
Fig.~\ref{fig:communication} provides continuous and consistent recording of
the analog parameters.

% -------------------------------------------------
% Figure 12: Light and soil-moisture time series
% -------------------------------------------------
\begin{figure*}[!t]
    \centering
    \includegraphics[width=\textwidth,height=6.8cm]{figA5.jpg}
    \caption{Light level (LDR) and soil-moisture analog readings over the
    five-day laboratory deployment. Both channels were multiplexed through
    the CD4051 and acquired using the single ADC channel of the ESP8266.}
    \label{fig:appendix_ml3}
\end{figure*}

\subsection{24-h Environmental Forecasting}

Figure~\ref{fig:predictions} presents the ``All Parameters--Predictions''
view of the dashboard, showing the predicted values of the five environmental
parameters over the subsequent 24~h from the time of data capture. The
horizontal axis represents the prediction time, while the vertical axis
represents normalized measurements in the range of 0--100, allowing the
different parameters to be compared within a common scale. The dashboard
also displays the predicted mean, minimum, and maximum values for each
parameter.

\begin{itemize}
    \item \textbf{Temperature:} A decreasing tendency is observed over the
    24-h prediction period, with the normalized measurement decreasing from
    approximately 100 to 0. The predicted temperature ranges from about
    16.7 to 21.6~$^\circ$C, with an average value of approximately
    19.1~$^\circ$C.

    \item \textbf{Humidity:} An increasing trend is observed, with the
    normalized measurement increasing from approximately 0 to 100. The
    predicted relative humidity ranges from about 53.6\% to 67.4\% RH,
    with an average value of approximately 60.5\%.

    \item \textbf{Gas:} The gas level remains relatively constant throughout
    the 24-h prediction period, with normalized fluctuations occurring
    approximately within the range of 50--60. The predicted gas level varies
    only slightly, with a minimum of approximately 77.0 and a maximum of
    approximately 77.6.

    \item \textbf{Light:} The light parameter exhibits a general downward
    trend over the 24-h period, which is consistent with the observed
    relationship between light intensity and temperature. The predicted
    light values range from approximately 314 to 532, with a mean value of
    approximately 423.

    \item \textbf{Soil moisture:} The normalized soil-moisture measurement
    exhibits an increasing trend similar to that of humidity. The reported
    statistics indicate a mean, minimum, and maximum value of approximately
    500. Therefore, the appropriateness of the original-unit representation
    should be considered when interpreting the soil-moisture prediction.
\end{itemize}

Overall, the 24-h predictions indicate a transition toward a cooler and more
humid environmental condition. Temperature and light intensity are predicted
to decrease, whereas humidity and soil moisture increase. The gas level
remains nearly constant, indicating relatively stable gas conditions during
the prediction period.

% -------------------------------------------------
% Figure 8
% -------------------------------------------------

\begin{figure}[!t]
    \centering
    \includegraphics[width=\columnwidth]{fig8.png}
    \caption{Twenty-four-hour predictions of temperature, humidity, gas level,
    light intensity, and soil moisture generated by the proposed EMD
    forecasting system.}
    \label{fig:predictions}
\end{figure}

% -------------------------------------------------
% Measurement Accuracy Validation
% -------------------------------------------------

\subsection{Measurement Accuracy Validation}

In order to confirm the overall measurement accuracy of the proposed
monitoring system, the sensor readings of the monitoring system were compared
with those from existing laboratory hardware that acts as the reference.
A total of 100 measurements were randomly selected and conducted manually
from the reference device at a certain point in time while the recorded
observations were obtained simultaneously by the proposed EMD. The date,
time, readings of the instruments, and corresponding readings of the developed
system are presented in Table~\ref{tab:measurement_verification}.

% -------------------------------------------------
% Table I
% -------------------------------------------------

\begin{table*}[!t]
    \centering
    \caption{Representative measurements from the laboratory reference
    equipment and the proposed EMD.}
    \label{tab:measurement_verification}

    \renewcommand{\arraystretch}{1.05}
    \setlength{\tabcolsep}{3.2pt}

    \begin{tabular}{c l cc cc cc cc cc}
        \toprule
        \multirow{2}{*}{No.} &
        \multirow{2}{*}{Date \& Time} &
        \multicolumn{2}{c}{Temperature ($^\circ$C)} &
        \multicolumn{2}{c}{Humidity (\%)} &
        \multicolumn{2}{c}{Gas} &
        \multicolumn{2}{c}{Light} &
        \multicolumn{2}{c}{Soil} \\
        
        \cmidrule(lr){3-4}
        \cmidrule(lr){5-6}
        \cmidrule(lr){7-8}
        \cmidrule(lr){9-10}
        \cmidrule(lr){11-12}
        
        &
        &
        Ref. & EMD &
        Ref. & EMD &
        Ref. & EMD &
        Ref. & EMD &
        Ref. & EMD \\
        
        \midrule
        
        1 & 11-Jan-2026 10:05 AM &
        22.4 & 22.5 & 50.2 & 50.1 & 777 & 785 & 205 & 185 & 500 & 500 \\

        2 & 11-Jan-2026 10:15 AM &
        22.7 & 22.6 & 49.8 & 49.9 & 787 & 795 & 375 & 375 & 500 & 500 \\

        3 & 11-Jan-2026 10:25 AM &
        23.1 & 23.0 & 49.3 & 49.4 & 807 & 815 & 485 & 505 & 500 & 500 \\

        4 & 11-Jan-2026 10:35 AM &
        23.5 & 23.6 & 48.9 & 49.6 & 818 & 826 & 605 & 585 & 500 & 500 \\

        5 & 11-Jan-2026 10:45 AM &
        23.8 & 23.7 & 48.5 & 48.6 & 838 & 845 & 755 & 745 & 500 & 500 \\

        6 & 11-Jan-2026 11:00 AM &
        24.1 & 24.0 & 47.9 & 48.0 & 858 & 865 & 590 & 592 & 499 & 500 \\

        7 & 11-Jan-2026 11:15 AM &
        24.3 & 24.4 & 47.5 & 47.6 & 878 & 885 & 605 & 603 & 499 & 500 \\

        8 & 11-Jan-2026 11:30 AM &
        24.0 & 24.1 & 47.8 & 47.7 & 848 & 855 & 590 & 591 & 498 & 500 \\

        9 & 12-Jan-2026 11:00 AM &
        23.2 & 23.3 & 49.1 & 49.0 & 797 & 805 & 405 & 425 & 498 & 500 \\

        10 & 12-Jan-2026 11:30 AM &
        22.8 & 22.7 & 50.0 & 50.1 & 767 & 775 & 510 & 509 & 500 & 500 \\

        \bottomrule
    \end{tabular}
\end{table*}

Table~\ref{tab:measurement_verification} presents 10 representative
observations out of the 100 observations used for verification. Ref. denotes
readings from the laboratory reference equipment, while EMD denotes readings
recorded by the proposed Environmental Monitoring Device.

The readings from both systems were matched and compared according to the
different measured parameters. The results showed that the temperature
values recorded by the EMD were approximately 0.1--0.2~$^\circ$C lower than
those obtained from the reference equipment. Humidity values differed by
approximately 0.1--0.7\%. Gas readings differed by approximately 7--8 units,
whereas light readings showed differences ranging from 0--20 units depending
on the measured intensity. Soil-moisture readings exhibited very small
deviations of 0--2 units, with several observations showing identical values.

The verification results showed a 94\% overall agreement between the proposed
system and the reference measurements across the evaluated observations.
The agreement was calculated as the ratio of the number of measurements
satisfying the predefined acceptance criterion to the total number of
evaluated observations, expressed as

\begin{equation}
    \mathrm{Agreement}~(\%) =
    \frac{N_{\mathrm{agreement}}}
    {N_{\mathrm{total}}}
    \times 100
    =
    \frac{94}{100}\times100
    =94\%.
    \label{eq:agreement}
\end{equation}

where $N_{\mathrm{agreement}}$ represents the number of measurements in
agreement with the reference equipment and $N_{\mathrm{total}}$ represents
the total number of evaluated observations.

The close correspondence between the measurements and their time-series
trends provides experimental evidence for the reliability of the proposed
sensing and acquisition system. The comparison was further evaluated using
corresponding time-series plots to examine the agreement and consistency
between the proposed EMD and the laboratory reference measurements.

\begin{table*}[!h]
    \centering
    \caption{Comparative analysis of the proposed system with other recent
    works in IoT environmental monitoring and machine learning for sensor data.}
    \label{tab:comparative_analysis}

    \scriptsize
    \renewcommand{\arraystretch}{1.15}
    \setlength{\tabcolsep}{3.5pt}

    \begin{tabular}{c c p{4.0cm} c c c c c c c}
        \toprule
        \textbf{Year} &
        \textbf{Ref.} &
        \textbf{Contribution} &
        \multicolumn{6}{c}{\textbf{Scope}} \\

        \cmidrule(lr){4-9}

        & &
        &
        \textbf{IoT env.} &
        \textbf{ADC/ Mux} &
        \textbf{Cloud/ back.} &
        \textbf{ML pred.} &
        \textbf{Anom. det.} &
        \textbf{Lab data} \\

        \midrule

        2019 & [1] &
        Low-cost IoT, MQTT, cloud air quality &
        $\checkmark$ & $\times$ & $\checkmark$ &
        $\times$ & $\times$ & $\times$ \\

        2021 & [2] &
        NodeMCU climate, Firebase; offline limits &
        $\checkmark$ & $\times$ & $\checkmark$ &
        $\times$ & $\times$ & $\times$ \\

        2022 & [7] &
        ESP8266 environmental data acquisition &
        $\checkmark$ & $\times$ & $\checkmark$ &
        $\times$ & $\times$ & $\times$ \\

        2023 & [8] &
        Gas detection, MQ sensors, ESP8266 &
        $\checkmark$ & $\times$ & $\checkmark$ &
        $\times$ & $\times$ & $\times$ \\

        2018 & [10] &
        Zigbee WSN environmental monitoring &
        $\checkmark$ & $\times$ & $\checkmark$ &
        $\times$ & $\times$ & $\times$ \\

        2020 & [11] &
        LoRaWAN scalable environmental monitoring &
        $\checkmark$ & $\times$ & $\checkmark$ &
        $\times$ & $\times$ & $\times$ \\

        2022 & [3] &
        ML anomaly detection in IoT data &
        $\checkmark$ & $\times$ & $\checkmark$ &
        $\checkmark$ & $\checkmark$ & $\times$ \\

        2022 & [21] &
        Modular sensor fusion; no CD4051 detail &
        $\checkmark$ & $\times$ & $\checkmark$ &
        $\times$ & $\times$ & $\times$ \\

        -- & \textbf{Proposed} &
        CD4051--ESP8266, 5 parameters, Flask/React,
        LSTM + Isolation Forest, lab data only &
        $\checkmark$ & $\checkmark$ & $\checkmark$ &
        $\checkmark$ & $\checkmark$ & $\checkmark$ \\

        \bottomrule
    \end{tabular}
\end{table*}

Table~\ref{tab:comparative_analysis} presents a comparison highlighting the
major differences between existing approaches in environmental monitoring
and the proposed system. Most of the reviewed studies have emphasized
IoT-based environmental monitoring or cloud integration, while some have
also incorporated machine-learning methods. The comparison indicates that
the combination of ADC/multiplexer-based acquisition, multi-parameter
monitoring, machine-learning-based prediction, anomaly detection, and
laboratory evaluation is less commonly represented within the compared
works.
% =================================================
% IV. DISCUSSION
% =================================================

\section{Discussion}

The proposed environmental monitoring system demonstrates the feasibility of
integrating multi-parameter sensing, wireless communication, data management,
visualization, and machine learning on a resource-constrained IoT platform.
Wireless sensor networks have previously demonstrated their effectiveness for
environmental monitoring applications~\cite{ref13}, while low-cost sensing
platforms have established the practicality of affordable environmental data
acquisition~\cite{ref14}. Low-cost sensing has also been investigated for
air-quality and exposure monitoring~\cite{ref15}. Edge-based IoT architectures
have further demonstrated the suitability of resource-constrained devices for
environmental monitoring~\cite{ref16}. Cloud-connected wireless sensor networks
provide an effective mechanism for transmitting and managing environmental
observations~\cite{ref17}. IoT-based smart-farming systems have shown the value
of continuous monitoring of environmental and agricultural parameters~\cite{ref18},
while irrigation monitoring systems demonstrate the application of
multi-parameter sensing in agricultural environments~\cite{ref19}.

The proposed system extends these approaches by combining temperature,
humidity, gas, light, and soil-moisture sensing within a single monitoring
node. Multi-sensor architectures have previously been applied to water-quality
monitoring using heterogeneous sensing modalities~\cite{ref20}. Systematic
studies of IoT-based environmental monitoring further highlight the importance
of integrating multiple sensing parameters within connected monitoring
architectures~\cite{ref21}. AI-based environmental monitoring has also
demonstrated the potential for combining real-time sensing with predictive
analytics~\cite{ref22}. In the present work, the CD4051 analog multiplexer
addresses the ADC limitation of the ESP8266 by allowing multiple analog
sensors to share a single ADC input. This hardware approach enables
heterogeneous sensors to be integrated without requiring a controller with
multiple ADC channels.

The backend architecture provides the connection between the sensing node and
user-level monitoring. Edge and IoT architectures have demonstrated the
importance of scalable processing for real-time monitoring applications~\cite{ref23}.
REST-based and cloud-oriented architectures provide mechanisms for collecting,
storing, and managing sensor observations~\cite{ref24}. The proposed
implementation follows this general approach using Flask for backend
communication, database storage for recorded measurements, and a React-based
dashboard for visualization. The resulting architecture allows measurements
to be acquired from the monitoring node, transmitted to the backend, stored,
and subsequently accessed for analysis~\cite{ref25}. Interactive visualization
further supports both real-time observation and retrospective analysis of
environmental measurements~\cite{ref26}. However, because the proposed system
was evaluated in a controlled laboratory environment, its performance should
not be interpreted as validation for direct agricultural or industrial
deployment.

Machine learning adds predictive and analytical capability beyond conventional
visualization. LSTM networks have been widely established for time-series
prediction because they can capture temporal dependencies within sequential
observations~\cite{ref27}. Previous environmental forecasting studies have
applied LSTM models to predict temperature and humidity variations~\cite{ref28}.
Hybrid machine-learning approaches have also been investigated for
environmental prediction and proactive climate management~\cite{ref29}. In
the proposed system, LSTM is used to forecast environmental parameters from
the recorded sensor measurements.

Anomaly detection provides a complementary function by identifying observations
that deviate from normal environmental patterns. Isolation Forest is an
established unsupervised method for detecting anomalous observations~\cite{ref30}.
Its application to environmental monitoring has demonstrated its suitability
for identifying unusual patterns in environmental datasets~\cite{ref31}.
Improved Isolation Forest approaches have also been investigated for anomaly
detection in wireless sensor networks~\cite{ref32}. Other IoT studies have
demonstrated the application of anomaly detection for identifying abnormal
system behavior and improving monitoring reliability~\cite{ref33}. More
recent work has combined unsupervised learning with IoT monitoring to
characterize operational states in laboratory environments~\cite{ref34}. In
this study, Isolation Forest is integrated with the LSTM forecasting pipeline
so that prediction and anomaly identification are performed within the same
monitoring framework rather than as independent post-processing tasks.

Using the information obtained from the implemented monitoring system makes
it possible to evaluate the machine-learning pipeline using recorded
measurements. However, evidence collected over five days in the laboratory
was limited by the variability of the environment and thus restricts the
generalizability of the models to different locations, seasons, and modes of
operation. Therefore, long-term data are necessary for further validation of
the models. Future work should investigate the integration of additional
sensors and monitoring nodes to further evaluate the effectiveness of the
multiplexed acquisition system.

% -------------------------------------------------
% Table II: Comparative Analysis
% -------------------------------------------------



In Table~\ref{tab:comparative_analysis}, a check mark represents a feature or
scope considered by a study, whereas a cross mark represents a feature not
identified in the compared work. The proposed framework incorporates
hardware-efficient data acquisition using a CD4051 multiplexer with an
ESP8266, monitoring of five environmental parameters, cloud/backend support,
LSTM-based prediction, Isolation Forest-based anomaly detection, and
laboratory data evaluation. Thus, the comparison demonstrates the broader
integrated scope of the proposed framework across sensing, communication,
data management, predictive analytics, anomaly detection, and experimental
evaluation.

% =================================================
% ABBREVIATIONS
% =================================================


%\clearpage

% =================================================
% V. CONCLUSION
% =================================================

\section{Conclusion}

The proposed Environmental Monitoring Device (EMD) represents a unified
solution for multi-parameter environmental monitoring while addressing the
limited availability of ADC channels in affordable IoT platforms. The
combination of the CD4051 analog multiplexer and the ESP8266 enables the
acquisition of environmental parameters, including temperature, humidity,
gas concentration, incident light intensity, and soil moisture, through a
single ADC channel. The system integrates a sensing unit with wireless
communication, a Flask-based backend, cloud data storage, a React-based
visualization interface, LSTM-based forecasting, and Isolation Forest-based
anomaly detection, thereby providing an integrated framework for environmental
data collection and analysis.

The system was modelled and tested for five consecutive days with a sampling
interval of 10~s. The experimental results demonstrated reliable operation
without noticeable channel-to-channel interference during multiplexed
acquisition. The machine-learning pipeline provided anomaly-detection
performance of approximately 92\% and enabled forecasting of the monitored
environmental parameters.

A key feature of the proposed system is the integration of hardware-efficient
measurement techniques with software-based data analysis within a single
platform. The availability of laboratory-measured data establishes a direct
link between the physical measurement system and subsequent analytical
evaluation. Nevertheless, the analysis is based on a five-day laboratory
dataset, which limits the generalizability of the findings across different
geographical locations, seasons, and longer-term environmental conditions.
Therefore, further evaluation using longer-duration measurements and diverse
environmental datasets is required.

Future investigations will focus on improving the forecasting and anomaly
detection capabilities of the proposed system and evaluating its scalability
across multiple monitoring locations. Lightweight transformer architectures,
including Temporal Fusion Transformer variants, may be investigated for
long-term forecasting. Hybrid approaches combining Isolation Forest with
sequence-based algorithms may also be explored to improve anomaly detection.
In addition, federated and adaptive online learning methods could be considered
to improve model adaptability across distributed EMD nodes.

% =================================================
% REFERENCES
% =================================================

%\clearpage

\bibliographystyle{IEEEtran}
\bibliography{references}

\end{document}
